\documentclass[11pt]{article}
\usepackage[margin=1in]{geometry}
\usepackage[T1]{fontenc}
\usepackage[utf8]{inputenc}
\usepackage{amsmath,amssymb,amsthm,mathtools}
\usepackage{enumitem}
\usepackage[hidelinks]{hyperref}

\newtheorem{theorem}{Theorem}
\newtheorem{lemma}{Lemma}
\newtheorem{corollary}{Corollary}
\newtheorem{proposition}{Proposition}
\newtheorem{remark}{Remark}
\newcommand{\R}{\mathbb{R}}

\newcommand{\e}{\mathrm{e}}
\newcommand{\N}{\mathbb{N}}
\newcommand{\diam}{\operatorname{diam}}
\newcommand{\Aut}{\operatorname{Aut}}
\newcommand{\conv}{\operatorname{conv}}
\newcommand{\tr}{\operatorname{tr}}
\newcommand{\relint}{\operatorname{relint}}
\DeclareMathOperator{\Log}{Log}

\title{Modernized proof notes: Lemmens--Nussbaum 2013, Birkhoff's version of Hilbert's metric}
\author{}
\date{2026-07-09}

\begin{document}
\maketitle

\section*{Source map and scope}
\begin{description}[style=nextline,leftmargin=3.2cm,labelwidth=3cm]
\item[Primary source] \texttt{\detokenize{Lemmens#U2013Nussbaum-2013-1304.7921v1.pdf}}
\item[Coverage] Sections 2--4: Hilbert's projective metric, order-cone formulas, nonexpansiveness and contraction, positive-matrix diameter, special-cone models, periodic-orbit results, and Denjoy--Wolff-type boundary dynamics.
\item[Source anchors]
\begin{itemize}[leftmargin=2em]
  \item Lemma 2.1
  \item Theorem 2.2
  \item Theorems 2.11 and 2.14
  \item Lemma 3.1
  \item Theorems 3.3--3.7
  \item Theorems 4.1 and 4.4--4.11, Corollary 4.2, and Remark 4.3
\end{itemize}
\end{description}

The exposition below preserves the source proof order and proof technique.  Notation is normalized
and intermediate calculations are made explicit.  When the source contains a gap, ambiguity, or
apparently incorrect step, the reconstruction records the source step before adding a separate
commentary remark.

\section*{Reference being modernized}
Bas Lemmens and Roger Nussbaum. ``Birkhoff's version of Hilbert's metric and its applications in analysis.'' arXiv:1304.7921v1, 2013.
The source is the 2013 survey by Lemmens and Nussbaum, arXiv:1304.7921.

\section*{Purpose of this note}
This reconstruction keeps the modern notation and the proof-level details supplied by the survey: the order definition of Hilbert's metric, the proof of its basic metric properties, the equivalence with Hilbert's cross-ratio metric, nonexpansiveness of order-preserving homogeneous maps, the Birkhoff contraction theorem as the main theorem seam, the finite positive-matrix diameter formula, the Perron--Frobenius operator application, and the special-cone formulas for simplicial, polyhedral, and symmetric cones.

The survey also contains later dynamical results for nonexpansive maps.  Those are stated with their source status rather than reproved when the survey itself cites an external proof.  The fully expanded material below covers the arguments actually supplied in the survey for Hilbert geometry, positive maps, and finite positive matrices.

\section{Cones, order, domination, and parts}
Let \(V\) be a real vector space.  A subset \(C\subseteq V\) is a cone if
\begin{enumerate}[label=(C\arabic*),leftmargin=2em]
\item \(C\) is convex;
\item \(\lambda C\subseteq C\) for all \(\lambda\ge0\);
\item \(C\cap(-C)=\{0\}\).
\end{enumerate}
If only the first two conditions hold, \(C\) is a wedge.  A cone induces a partial order
\[
  x\le_C y\quad\Longleftrightarrow\quad y-x\in C.
\]
A wedge induces only a preorder, because antisymmetry may fail.

If \(C\) is a closed cone in a Banach space \((V,\|\cdot\|)\), it is called normal if there exists \(\kappa>0\) such that
\[
  0\le_C x\le_C y\quad\Longrightarrow\quad \|x\|\le\kappa\|y\|.
\]
All closed cones in finite-dimensional normed spaces are normal, but infinite-dimensional closed cones need not be.

The dual wedge is
\[
  C^*=\{\varphi\in V^*: \varphi(x)\ge0\text{ for all }x\in C\}.
\]
If \(C-C\) is dense in \(V\), then \(C^*\) is a cone.  In finite dimensions, if \(C\) is closed and has nonempty interior, then \(C^*\) is a closed cone with nonempty interior.

For \(\varphi\in C^*\), define the base
\[
  \Sigma_\varphi=\{x\in C:\varphi(x)=1\}.
\]
In finite dimensions, if \(C\) is closed with nonempty interior and \(\varphi\in (C^*)^\circ\), then \(\Sigma_\varphi\) is nonempty, compact, and convex.  In infinite dimensions such bounded bases may fail to exist; the standard cone in \(\ell^2\) is one example.

\subsection{Domination}
For \(x\in V\) and \(y\in C\), say that \(y\) dominates \(x\) if there exist \(\alpha,\beta\in\R\) such that
\[
  \alpha y\le_C x\le_C \beta y.
\]
When \(y\) dominates \(x\), define
\[
  M(x/y)=\inf\{\beta\in\R:x\le_C \beta y\},
  \qquad
  m(x/y)=\sup\{\alpha\in\R:\alpha y\le_C x\}.
\]
If \(y=0\) and \(y\) dominates \(x\), then \(x=0\), because \(x\in C\cap(-C)\).  If \(y\ne0\), then \(M(x/y)\ge m(x/y)\).

Define an equivalence relation on \(C\) by
\[
  x\sim_C y
  \quad\Longleftrightarrow\quad
  x\text{ dominates }y\text{ and }y\text{ dominates }x.
\]
The equivalence classes are called parts of \(C\).  The origin is a part by itself.  If \(C\) is closed with nonempty interior in a Banach space, then \(C^\circ\) is a part.  For finite-dimensional closed cones, parts are precisely the relative interiors of faces.

For \(x,y\in C\setminus\{0\}\), comparability is equivalent to the existence of \(0<\alpha\le\beta\) with
\[
  \alpha y\le_C x\le_C \beta y.
\]
Also
\begin{equation}
  m(x/y)=\sup\{\alpha>0:y\le_C \alpha^{-1}x\}=M(y/x)^{-1}.
  \label{eq:LN-m-inverse-M}
\end{equation}

\section{Birkhoff's version of Hilbert's metric}
For comparable points \(x\sim_C y\) with \(y\ne0\), define
\begin{equation}
  d(x,y)=\log\frac{M(x/y)}{m(x/y)}.
  \label{eq:LN-hilbert-metric}
\end{equation}
Set \(d(0,0)=0\), and set \(d(x,y)=\infty\) when \(x\) and \(y\) are not comparable.

\begin{lemma}[Lemma 2.1]
Let \(C\) be a cone in \(V\).  If \(x\sim_C y\sim_C z\) with \(y\ne0\), then:
\begin{enumerate}[label=(\roman*),leftmargin=2em]
\item \(d(x,y)\ge0\);
\item \(d(x,y)=d(y,x)\);
\item \(d(x,z)\le d(x,y)+d(y,z)\);
\item \(d(\lambda x,\mu y)=d(x,y)\) for all \(\lambda,\mu>0\).
\end{enumerate}
If \(C\) is a closed cone in a Banach space, then \(d(x,y)=0\) if and only if \(x=\lambda y\) for some \(\lambda\ge0\).
\end{lemma}

\begin{proof}
For nonnegativity, take \(0<\alpha<m(x/y)\) and \(\beta>M(x/y)\).  Then
\[
  \alpha y\le_C x\le_C \beta y.
\]
Thus \((\beta-\alpha)y\in C\).  Since \(y\in C\setminus\{0\}\) and \(C\cap(-C)=\{0\}\), this forces \(\beta\ge\alpha\).  Letting \(\alpha\uparrow m(x/y)\) and \(\beta\downarrow M(x/y)\) yields \(M(x/y)/m(x/y)\ge1\), hence \(d(x,y)\ge0\).

Symmetry follows from \eqref{eq:LN-m-inverse-M}:
\[
  d(x,y)=\log\bigl(M(x/y)M(y/x)\bigr)=d(y,x).
\]
For the triangle inequality, let \(0<\alpha<m(x/y)\) and \(0<\gamma<m(y/z)\).  Then \(\alpha y\le_C x\) and \(\gamma z\le_C y\), hence \(\alpha\gamma z\le_C x\).  Thus
\[
  m(x/z)\ge m(x/y)m(y/z).
\]
Similarly, if \(x\le_C \beta y\) and \(y\le_C \delta z\), then \(x\le_C\beta\delta z\), so
\[
  M(x/z)\le M(x/y)M(y/z).
\]
Divide and take logarithms to obtain (iii).

For projective invariance, direct substitution gives
\[
  M(\lambda x/\mu y)=\frac{\lambda}{\mu}M(x/y),
  \qquad
  m(\lambda x/\mu y)=\frac{\lambda}{\mu}m(x/y),
\]
so the quotient in \eqref{eq:LN-hilbert-metric} is unchanged.

Finally assume \(C\) is closed and \(x\sim_C y\).  Closedness implies the endpoint inequalities are attained:
\[
  m(x/y)y\le_C x\le_C M(x/y)y.
\]
If \(d(x,y)=0\), then \(M(x/y)=m(x/y)\).  Hence
\[
  x=m(x/y)y.
\]
Conversely, if \(x=\lambda y\) with \(\lambda>0\), projective invariance gives \(d(x,y)=0\).  If \(\lambda=0\), then \(x=0\), comparable only to \(0\), and the convention gives \(d(0,0)=0\).
\end{proof}

\section{Equivalence with Hilbert's cross-ratio metric}
Let \(C\) be a closed cone with nonempty interior in an \((n+1)\)-dimensional real normed space \(V\).  Let \(H\subset V\) be an \(n\)-dimensional affine hyperplane such that
\[
  \Omega_C=H\cap C^\circ
\]
is a relatively open bounded convex set in \(H\).

\begin{theorem}[Theorem 2.2]
The restriction of \(d\) to \(\Omega_C\) equals Hilbert's cross-ratio metric \(\delta\) on \(\Omega_C\).
\end{theorem}

\begin{proof}
Take distinct \(x,y\in\Omega_C\).  Put
\[
  \alpha=m(x/y)=M(y/x)^{-1},
  \qquad
  \beta=M(x/y).
\]
Closedness gives
\[
  u=x-\alpha y\in\partial C,
  \qquad
  w=y-x/\beta\in\partial C.
\]
Let \(\ell_{xy}\) be the affine line through \(x\) and \(y\).  Let \(x'\) and \(y'\) be the two boundary points in \(\partial\Omega_C\) ordered so that \(x\) lies between \(x'\) and \(y\), and \(y\) lies between \(x\) and \(y'\).  Then there exist parameters \(\sigma,\tau>1\) such that
\[
  x'=y+\sigma(x-y),
  \qquad
  y'=x+\tau(y-x).
\]
Let \(\varphi\) be an affine-normalizing linear functional with
\[
  H=\{z\in V:\varphi(z)=1\}.
\]
Since \(x,y\in H\), \(\varphi(u)=1-\alpha\), and the ray through \(u\) meets \(H\) at
\[
  \frac{u}{\varphi(u)}=\frac{x-\alpha y}{1-\alpha}.
\]
Equating this with \(x'=y+\sigma(x-y)=\sigma x+(1-
\sigma)y\) gives
\[
  \frac1{1-\alpha}=\sigma,
  \qquad
  -\frac{\alpha}{1-\alpha}=1-\sigma.
\]
Thus \(\alpha=(\sigma-1)/\sigma\).  Similarly, using \(w=y-x/\beta\) and \(\varphi(w)=1-1/\beta\),
\[
  y'=\frac{w}{\varphi(w)}=\frac{y-x/\beta}{1-1/\beta}
\]
gives \(\beta=\tau/(\tau-1)\), up to the usual orientation convention.  Therefore the ratios of Euclidean lengths along the line satisfy
\[
  \frac{\|y-x'\|}{\|x-x'\|}=\frac{1}{\alpha}=M(y/x),
  \qquad
  \frac{\|x-y'\|}{\|y-y'\|}=\beta=M(x/y).
\]
Hilbert's cross-ratio formula is
\[
  \delta(x,y)=\log\left(
  \frac{\|y-x'\|}{\|x-x'\|}
  \frac{\|x-y'\|}{\|y-y'\|}
  \right).
\]
Substituting the preceding identities yields
\[
  \delta(x,y)=\log(M(y/x)M(x/y))
  =\log\frac{M(x/y)}{m(x/y)}=d(x,y),
\]
using \eqref{eq:LN-m-inverse-M}.
\end{proof}

\begin{remark}[Funk metric]
The asymmetric quantity
\[
  F_C(x,y)=\log M(x/y)
\]
is the Funk weak metric.  Hilbert's metric is the symmetrization
\[
  d(x,y)=F_C(x,y)+F_C(y,x).
\]
\end{remark}

\begin{remark}[Open convex sets in Banach spaces]
If \(\Omega\subset V\) is open, convex, and contains no affine line, set
\[
  C_\Omega=\{(s,sv):s>0,\;v\in\Omega\}\subset \R\times V.
\]
The closure \(\overline{C_\Omega}\) is a cone.  If \(w=(0,u)\) and \(-w\) both belonged to the closure, then by approximating \(w\) and \(-w\) using points \((s_k,s_kv_k)\), convexity would imply \(v+tu\in\Omega\) for all \(t\in\R\), contradicting the no-line assumption.  Thus the cone construction gives Hilbert's metric on arbitrary open line-free convex sets by restricting to the slice \(s=1\).
\end{remark}

\section{Order-preserving homogeneous maps are nonexpansive}
Let \(C\subset V\), \(K\subset W\) be cones.  A map \(f:C\to K\) is order-preserving if
\[
  x\le_C y\quad\Longrightarrow\quad f(x)\le_K f(y),
\]
and homogeneous of degree \(r\) if \(f(\lambda x)=\lambda^r f(x)\) for all \(\lambda>0\).

\begin{proposition}[Proposition 2.6]
If \(f:C\to W\) is order-preserving and homogeneous of degree \(r>0\), then for all \(x\sim_C y\),
\[
  M(f(x)/f(y))\le M(x/y)^r,
  \qquad
  m(x/y)^r\le m(f(x)/f(y)).
\]
\end{proposition}

\begin{proof}
If \(x\le_C \beta y\), then order preservation gives
\[
  f(x)\le_K f(\beta y)=\beta^r f(y).
\]
Taking the infimum over \(\beta>M(x/y)\) yields
\[
  M(f(x)/f(y))\le M(x/y)^r.
\]
Similarly, if \(\alpha y\le_C x\), then
\[
  \alpha^r f(y)=f(\alpha y)\le_K f(x),
\]
so \(m(f(x)/f(y))\ge \alpha^r\).  Taking the supremum over \(\alpha<m(x/y)\) gives the second inequality.
\end{proof}

\begin{corollary}[Corollary 2.7]
If \(f:C\to K\) is order-preserving and homogeneous of degree \(1\), then \(f\) is nonexpansive for Hilbert's metric on each part:
\[
  d_K(f(x),f(y))\le d_C(x,y).
\]
\end{corollary}

\begin{proof}
Use Proposition 2.6 with \(r=1\):
\[
  \frac{M(f(x)/f(y))}{m(f(x)/f(y))}
  \le \frac{M(x/y)}{m(x/y)}.
\]
Taking logarithms gives the result.
\end{proof}

Order-reversing homogeneous degree \(-1\) maps are also nonexpansive: the inequalities are reversed and inverted, but the quotient defining \(d\) is unchanged.  This is one reason the Hilbert metric is useful for alternating scaling maps, where one normalization can be order-reversing after reciprocation.

\section{Birkhoff contraction ratio and the finite positive matrix formula}
For a linear map \(L:V\to W\) with \(L(C)\subseteq K\), define
\[
  \kappa(L)=\inf\{c\ge0:d_K(Lx,Ly)\le c\,d_C(x,y)\text{ for all }x\sim_C y\},
\]
and
\[
  \Delta(L)=\sup\{d_K(Lx,Ly):x,y\in C,\;Lx\sim_K Ly\}.
\]

\begin{theorem}[Birkhoff, Theorem 2.9]
For a linear map \(L:V\to W\) with \(L(C)\subseteq K\),
\[
  \kappa(L)=\tanh(\Delta(L)/4),
\]
where \(\tanh(\infty)=1\).
\end{theorem}

This theorem is the same theorem modernized in the Eveson--Nussbaum and Birkhoff notes.  The survey uses it as a black-box contraction theorem.  For a strictly positive \(m\times n\) matrix \(A=(a_{pi})\), the survey records the explicit formula
\begin{equation}
  \Delta(A)=\max_{i,j} d(Ae_i,Ae_j)
  =\log\max_{i,j,p,q}\frac{a_{pi}a_{qj}}{a_{pj}a_{qi}}<\infty.
  \label{eq:positive-matrix-diameter}
\end{equation}

\subsection{Proof of the matrix diameter formula}
For \(x\in\R_+^n\setminus\{0\}\),
\[
  Ax=\sum_i x_i Ae_i.
\]
The image \(A(\R_+^n)\) is the cone generated by the columns \(Ae_i\).  Proposition 2.9 in the Eveson--Nussbaum notation, or the convex-hull invariance of projective diameter, gives that the diameter of the image cone is the diameter of its set of generating rays:
\[
  \Delta(A)=\max_{i,j}d(Ae_i,Ae_j).
\]
For two columns \(Ae_i=(a_{pi})_p\) and \(Ae_j=(a_{pj})_p\), the positive-orthant formula gives
\[
  d(Ae_i,Ae_j)=
  \log\frac{\max_p a_{pi}/a_{pj}}{\min_q a_{qi}/a_{qj}}.
\]
The reciprocal of the minimum is a maximum, so
\[
  \frac{\max_p a_{pi}/a_{pj}}{\min_q a_{qi}/a_{qj}}
  =\max_{p,q}\frac{a_{pi}/a_{pj}}{a_{qi}/a_{qj}}
  =\max_{p,q}\frac{a_{pi}a_{qj}}{a_{pj}a_{qi}}.
\]
Taking the maximum over \(i,j\) gives \eqref{eq:positive-matrix-diameter}.  Strict positivity is exactly what makes every denominator positive and the maximum finite.

\section{Cone-linear Perron--Frobenius theorem}
A map \(L:C\to C\) is cone-linear if
\[
  L(\alpha x+\beta y)=\alpha Lx+\beta Ly
\]
for all \(\alpha,\beta\ge0\) and \(x,y\in C\).  The survey quotes the following theorem.

\begin{theorem}[Theorem 2.11, quoted theorem seam]
Let \(C\) be a closed normal cone in a Banach space \(V\), and let \(L:C\to C\) be cone-linear and continuous at \(0\).  If there exists \(p\ge1\) such that \(\Delta(L^p)<\infty\) and \(L^{p+1}(C)\ne\{0\}\), then \(L\) has a unique eigenvector \(v\in C\), \(\|v\|=1\), satisfying
\[
  Lv=r_C(L)v,
\]
where
\[
  r_C(L)=\lim_{k\to\infty}\|L^k\|_C^{1/k},
  \qquad
  \|L^k\|_C=\sup\{\|L^k x\|:x\in C,\|x\|=1\}.
\]
Moreover, if \(c=\tanh(\Delta(L^p)/4)<1\), then
\[
  d(L^{kp}x,v)\le c^k d(x,v)
\]
for all nonzero \(x\in C\) in the relevant part.
\end{theorem}

For finite positive matrices, this theorem collapses to Perron--Frobenius: strict positivity gives finite diameter for \(L\), hence a unique positive eigen-ray and geometric convergence in Hilbert metric.

\section{Application to Perron--Frobenius operators}
Let \((S,\rho)\) be a bounded complete metric space with diameter \(\Delta>0\).  Let \(C_b(S)\) be bounded continuous real-valued functions with the sup norm.  For \(0<\lambda\le1\) and \(M>0\), define
\begin{equation}
  K(M,\lambda)=\{f\in C_b(S): f(s)\le f(t)\exp(M\rho(s,t)^\lambda)\text{ for all }s,t\in S\}.
  \label{eq:holder-cone}
\end{equation}
This is a closed normal cone.  If \(f\in K(M,\lambda)\), then either \(f\equiv0\) or \(f(s)>0\) for all \(s\), and
\begin{equation}
  \sup_s f(s)\le \exp(M\Delta^\lambda)\inf_s f(s).
  \label{eq:holder-cone-sup-inf}
\end{equation}

Fix functions \(b_i\in K(M_0,\lambda)\) and Lipschitz maps \(\theta_i:S\to S\) with
\[
  \operatorname{Lip}(\theta_i)\le c<1.
\]
Assume \(\sum_i b_i(s_*)<\infty\) for some \(s_*\).  Define the Perron--Frobenius operator
\begin{equation}
  (Lf)(t)=\sum_i b_i(t)f(\theta_i(t)).
  \label{eq:LN-PF-operator}
\end{equation}
Then \(L:C_b(S)\to C_b(S)\) is bounded, because
\[
  \|Lf\|\le \sum_i b_i(s_*)\exp(M_0\Delta^\lambda)\|f\|.
\]

\begin{lemma}[Lemma 2.12]
Let \(d_2\) be Hilbert's metric on \(K(M_2,\lambda)\).  If \(0<M_1<M_2\), then \(K(M_1,\lambda)\setminus\{0\}\) lies in a single part of \(K(M_2,\lambda)\), and has finite \(d_2\)-diameter.
\end{lemma}

\begin{proof}
Let \(\le_2\) be the order induced by \(K(M_2,\lambda)\).  It suffices, by projective invariance, to take \(f,g\in K(M_1,\lambda)\) with \(\|f\|=\|g\|=1\) and find uniform \(0<\alpha\le\beta\) such that
\[
  \alpha f\le_2 g\le_2 \beta f.
\]
By symmetry, it suffices to prove the lower inequality.  The inequality \(\alpha f\le_2 g\) means
\begin{equation}
  g(s)-\alpha f(s)
  \le (g(t)-\alpha f(t))\exp(M_2\rho(s,t)^\lambda)
  \label{eq:alpha-lower-holder}
\end{equation}
for all \(s,t\).

Since \(f,g\in K(M_1,\lambda)\) and have sup norm \(1\), \eqref{eq:holder-cone-sup-inf} gives
\[
  \exp(-M_1\Delta^\lambda)
  \le f(t),g(t)
  \le1,
\]
and hence
\begin{equation}
  \exp(-M_1\Delta^\lambda)
  \le \frac{g(t)}{f(t)}
  \le \exp(M_1\Delta^\lambda).
  \label{eq:ratio-holder-bound}
\end{equation}
Choose \(\alpha<\exp(-M_1\Delta^\lambda)\) so that all denominators below are positive.  Dividing \eqref{eq:alpha-lower-holder} by \(g(t)-\alpha f(t)\), it is enough to show
\[
  \frac{g(s)-\alpha f(s)}{g(t)-\alpha f(t)}
  \le \exp(M_2\rho(s,t)^\lambda).
\]
Using the \(K(M_1,\lambda)\) inequalities for \(g\) and \(f\),
\[
  g(s)\le g(t)\e^{M_1 r_0},
  \qquad
  f(s)\ge f(t)\e^{-M_1 r_0},
  \qquad r_0=\rho(s,t)^\lambda,
\]
so
\[
  \frac{g(s)-\alpha f(s)}{g(t)-\alpha f(t)}
  \le
  \frac{g(t)\,\e^{M_1r_0}-\alpha f(t)\e^{-M_1r_0}}{g(t)-\alpha f(t)}.
\]
Put \(r=M_2r_0\) and \(\mu=M_1/M_2<1\).  Using \eqref{eq:ratio-holder-bound}, the last ratio is bounded by \(\e^r\) provided
\begin{equation}
  \left(\frac{\alpha}{\e^{-M_1\Delta^\lambda}-\alpha}\right)
  \left(\frac{\e^{\mu r}-\e^{-\mu r}}{\e^r-\e^{\mu r}}\right)
  \le1.
  \label{eq:LN-theta-condition}
\end{equation}
The second factor extends continuously to \(r=0\) with value \(2\mu/(1-\mu)\), tends to \(0\) as \(r\to\infty\), and has negative derivative on \((0,\infty)\).  Thus its maximum is \(2\mu/(1-\mu)\).
Set
\[
  \kappa=\frac{M_2-M_1}{M_2+M_1}=\frac{1-\mu}{1+\mu},
  \qquad
  \alpha=\kappa\e^{-M_1\Delta^\lambda}.
\]
Then
\[
  \frac{\alpha}{\e^{-M_1\Delta^\lambda}-\alpha}
  =\frac{\kappa}{1-\kappa}
  =\frac{1-\mu}{2\mu},
\]
so \eqref{eq:LN-theta-condition} holds.  By symmetry one can take
\[
  \beta=\frac{M_2+M_1}{M_2-M_1}\e^{M_1\Delta^\lambda}.
\]
Consequently
\[
  d_2(f,g)\le \log(\beta/\alpha)
  =2\log\frac{M_2+M_1}{M_2-M_1}+2M_1\Delta^\lambda,
\]
a finite bound independent of \(f,g\).
\end{proof}

\begin{lemma}[Lemma 2.13]
Assume at least one \(b_i\) is not identically zero.  If \(M_2>M_0/(1-c^\lambda)\) and \(M_1=M_0+c^\lambda M_2\), then
\[
  L(K(M_2,\lambda)\setminus\{0\})\subseteq K(M_1,\lambda)\setminus\{0\}.
\]
\end{lemma}

\begin{proof}
For \(f\in K(M_2,\lambda)\), the composition \(f\circ\theta_i\) satisfies
\[
  f(\theta_i(s))
  \le f(\theta_i(t))
  \exp(M_2\rho(\theta_i(s),\theta_i(t))^\lambda)
  \le f(\theta_i(t))
  \exp(c^\lambda M_2\rho(s,t)^\lambda).
\]
Thus \(f\circ\theta_i\in K(c^\lambda M_2,\lambda)\).  Since \(b_i\in K(M_0,\lambda)\), multiplying inequalities gives
\[
  b_i(s)f(\theta_i(s))
  \le b_i(t)f(\theta_i(t))
  \exp((M_0+c^\lambda M_2)\rho(s,t)^\lambda).
\]
Hence each summand belongs to \(K(M_1,\lambda)\), and closedness of the cone gives \(Lf\in K(M_1,\lambda)\).

If \(f\ne0\), then \(f>0\) everywhere.  At least one \(b_j\) is not zero, and a nonzero element of \(K(M_0,\lambda)\) is positive everywhere.  Thus
\[
  (Lf)(t)\ge b_j(t)f(\theta_j(t))>0
\]
for all \(t\), so \(Lf\ne0\).
\end{proof}

\begin{theorem}[Theorem 2.14]
Under the preceding assumptions, for every \(M_2>M_0/(1-c^\lambda)\) there exists a unique \(v\in K(M_2,\lambda)\), \(\|v\|=1\), such that
\[
  Lv=\|Lv\|v.
\]
For every \(g\in K(M_2,\lambda)\setminus\{0\}\),
\[
  \lim_{k\to\infty}\left\|\frac{L^k g}{\|L^k g\|}-v\right\|=0.
\]
\end{theorem}

\begin{proof}
Let \(L_0\) be \(L\) restricted to \(K(M_2,\lambda)\).  By Lemma 2.13, \(L_0\) maps nonzero elements into \(K(M_1,\lambda)\setminus\{0\}\), where \(M_1=M_0+c^\lambda M_2<M_2\).  By Lemma 2.12, \(K(M_1,\lambda)\setminus\{0\}\) has finite diameter in the metric of \(K(M_2,\lambda)\).  Hence \(\Delta(L_0)<\infty\).  Also \(L_0^2(K(M_2,\lambda))\ne\{0\}\).  Apply Theorem 2.11.
\end{proof}

\section{Special cones}

\subsection{Extreme functionals and a general formula}
Let \(C\) be a closed cone with nonempty interior in a finite-dimensional vector space \(V\).  For \(u\in C^\circ\), define
\[
  \Sigma_u^*=\{\varphi\in C^*: \varphi(u)=1\}.
\]
This is compact convex.  Let \(E_u\) be the closure of its extreme points.  By Krein--Milman/Minkowski, \(\Sigma_u^*\) is the closed convex hull of its extreme points.

\begin{lemma}[Lemma 3.1]
For \(x\in V\) and \(y\in C^\circ\),
\[
  M(x/y)=\max_{\varphi\in E_u}\frac{\varphi(x)}{\varphi(y)},
  \qquad
  m(x/y)=\min_{\varphi\in E_u}\frac{\varphi(x)}{\varphi(y)}.
\]
\end{lemma}

\begin{proof}
Order is detected by the dual cone: \(x\le_C y\) if and only if \(\varphi(x)\le\varphi(y)\) for all \(\varphi\in C^*\).  If \(y-x\notin C\), Hahn--Banach separation gives a functional nonnegative on \(C\) but negative on \(y-x\), contradicting the dual inequalities.  Since every element of \(\Sigma_u^*\) is a convex combination/limit of extreme functionals, it suffices to test on \(E_u\).

Now \(x\le_C \beta y\) is equivalent to
\[
  \varphi(x)\le \beta\varphi(y)
\]
for all \(\varphi\in E_u\).  Since \(y\in C^\circ\), \(\varphi(y)>0\).  Therefore the least admissible \(\beta\) is the displayed maximum.  The proof for \(m\) is the same with the inequality \(\alpha y\le_C x\).
\end{proof}

\subsection{Simplicial cones}
For the standard positive cone \(\R_+^n\),
\[
  M(x/y)=\max_i\frac{x_i}{y_i},
  \qquad
  m(x/y)=\min_i\frac{x_i}{y_i}.
\]
The logarithmic coordinate map
\[
  \Log(x)=(\log x_1,\ldots,\log x_n)
\]
is an isometry from \(((\R_+^n)^\circ,d)\) to \(\R^n\) equipped with the variation seminorm
\[
  \|w\|_{\mathrm{var}}=\max_i w_i-\min_i w_i.
\]
Indeed,
\[
  d(x,y)=\max_i(\log x_i-\log y_i)-\min_j(\log x_j-\log y_j).
\]
On a cross-section, e.g. \(x_n=1\), this becomes a genuine norm on \(\R^{n-1}\):
\[
  \|z\|_H=\max\{z_1,
\ldots,z_{n-1},0\}-\min\{z_1,
\ldots,z_{n-1},0\}.
\]
Thus the Hilbert geometry on an open simplex is isometric to this normed vector space.

\subsection{Polyhedral cones}
If \(C\subset\R^n\) is a polyhedral cone with nonempty interior and facets defined by normalized functionals \(\psi_1,
\ldots,\psi_m\), then \(E_u=\{\psi_1,
\ldots,\psi_m\}\).  For \(x,y\in C^\circ\), Lemma 3.1 gives
\[
  d(x,y)=\log\max_{i,j}\frac{\psi_i(x)\psi_j(y)}{\psi_i(y)\psi_j(x)}.
\]
Define
\[
  \Psi_{ij}(x)=\log\frac{\psi_i(x)}{\psi_j(x)},
  \qquad 1\le i<j\le m.
\]
Then
\[
  d(x,y)=\max_{1\le i<j\le m}|\Psi_{ij}(x)-\Psi_{ij}(y)|.
\]
Hence the Hilbert geometry on the interior of a polytope with \(m\) facets embeds isometrically into
\[
  (\R^{m(m-1)/2},\|\cdot\|_\infty).
\]

\subsection{Symmetric cones}
Let \(V\) be a Euclidean Jordan algebra with symmetric cone \(C^\circ\).  For \(x\in C^\circ\), the quadratic representation \(P(x)\) is an automorphism of \(C\), and \(P(x^{-1/2})x=e\).  For \(w\in V\), define
\[
  \lambda_+(w,x)=\lambda_+(P(x^{-1/2})w),
  \qquad
  \lambda_-(w,x)=\lambda_-(P(x^{-1/2})w),
\]
where \(\lambda_+\) and \(\lambda_-\) are the largest and smallest Jordan eigenvalues.

\begin{theorem}[Theorem 3.7]
For \(w\in V\) and \(x\in C^\circ\),
\[
  M(w/x)=\lambda_+(w,x),
  \qquad
  m(w/x)=\lambda_-(w,x).
\]
Consequently
\[
  d(x,y)=\log\frac{\lambda_+(x,y)}{\lambda_-(x,y)}.
\]
\end{theorem}

\begin{proof}
Let \(z=P(x^{-1/2})w\), and write its spectral decomposition as
\[
  z=\lambda_1c_1+
\cdots+\lambda_kc_k,
\]
where \(c_i\) are a complete system of orthogonal idempotents.  Since \(P(x^{-1/2})\in\Aut(C)\),
\[
  w\le_C \beta x
  \quad\Longleftrightarrow\quad
  z\le_C \beta e.
\]
Because \(e=c_1+
\cdots+c_k\), this is equivalent to
\[
  \sum_i(\beta-
\lambda_i)c_i\in C,
\]
which holds exactly when \(\beta\ge\max_i\lambda_i\).  Hence \(M(w/x)=\lambda_+(w,x)\).  The proof for \(m\) is identical with \(\alpha x\le_C w\), giving \(\alpha\le\min_i\lambda_i\).
\end{proof}

For positive definite matrices, this specializes to
\[
  \lambda_+(A,B)=\lambda_{\max}(B^{-1/2}AB^{-1/2}),
  \qquad
  \lambda_-(A,B)=\lambda_{\min}(B^{-1/2}AB^{-1/2}).
\]

\section{Nonexpansive dynamics on Hilbert geometries}
The survey's Section 4 changes status: most results are quoted from the cited literature rather than
proved in the survey.  A faithful reconstruction must therefore preserve both the statements and their
source status.  The deductions made inside the survey---notably Corollary 4.2 from Theorems 3.5 and
4.1---are written out below; externally cited theorems are not presented as arguments supplied by
Lemmens and Nussbaum.

Let \((\Omega,\delta)\) be a Hilbert geometry.  Important nonlinear examples arise from a
continuous order-preserving homogeneous map \(f:C^\circ\to C^\circ\) and its normalization on a
base
\[
 \Sigma_\varphi^\circ=\{x\in C^\circ:\varphi(x)=1\},
 \qquad
 g(x)=\frac{f(x)}{\varphi(f(x))}.
\]
By Proposition 2.6 and Corollary 2.7, both maps are Hilbert-nonexpansive on each part.  Section 4
uses this fact as the bridge from cone order to metric dynamics.

\subsection{Periodic orbits when a fixed point exists}
For a map \(f\), write
\[
 \omega(x;f)=\{y:\ f^{k_i}(x)\to y\text{ for some }k_i\to\infty\}.
\]
An orbit converges to a periodic orbit of period \(p\) when \(f^{kp}(x)\) converges to a periodic
point of least period \(p\).

\begin{theorem}[Theorem 4.1, quoted from the finite-dimensional sup-norm theory]
Let \(X\) be a closed subset of \((\R^N,\|\cdot\|_\infty)\), and let
\(f:X\to X\) be nonexpansive with a fixed point.  Every orbit converges to a periodic orbit.  Its
period is at most
\[
 \max_{0\le k\le N}2^k\binom Nk.
\]
\end{theorem}
The survey cites the proof and historical discussion to the nonlinear Perron--Frobenius literature;
the numerical bound is due to Lemmens and Scheutzow.  No proof is supplied in the survey.

\begin{corollary}[Corollary 4.2]
If \((\Omega,\delta)\) is a polytopal Hilbert geometry whose defining polytope has \(m\) facets,
and \(f:D\to D\) is a nonexpansive self-map of a closed subset with a fixed point, then every orbit
converges to a periodic orbit of period at most
\[
 \max_{0\le k\le N}2^k\binom Nk,
 \qquad N=\frac{m(m-1)}2.
\]
\end{corollary}
\begin{proof}
Theorem 3.5 embeds \((\Omega,\delta)\) isometrically into
\((\R^N,\|\cdot\|_\infty)\) using the logarithmic facet-ratio coordinates
\(\Psi_{ij}(x)=\log(\psi_i(x)/\psi_j(x))\).  Conjugate \(f\) by this embedding and apply Theorem
4.1 to the closed embedded domain.  Pull the periodic limit back through the isometry.
\end{proof}

Remark 4.3 observes a geometric consequence: a Euclidean plane cannot embed isometrically in a
polytopal Hilbert geometry, since an irrational rotation would be a nonexpansive map with bounded
orbits that do not converge to periodic orbits.  Thus the Euclidean rank of a polytopal Hilbert
geometry is one.

The survey next records the sharper cone-specific classification.
\begin{theorem}[Theorem 4.4, quoted]
Let \(C\) be a polyhedral cone with nonempty interior and \(m\) facets, and let
\(f:C\to C\) be continuous, order-preserving, and homogeneous.
\begin{enumerate}[leftmargin=2em]
\item Every norm-bounded orbit converges to a periodic orbit, with period at most
\[
 \frac{m!}{\lfloor m/3\rfloor!\,\lfloor(m+1)/3\rfloor!\,\lfloor(m+2)/3\rfloor!}.
\]
\item If \(C\) is an \(n\)-dimensional simplicial cone, the possible periods are exactly the
integers \(p=q_1q_2\) for which, for some \(0\le k\le n\),
\[
 1\le q_1\le\binom{k}{\lfloor k/2\rfloor},
 \qquad
 1\le q_2\le\binom{n}{k}.
\]
\item If \(f(v)=\lambda v\) for some \(v\in C^\circ\), then every orbit of the normalized map
\(g\) converges to a periodic orbit of period at most \(\binom{m}{\lfloor m/2\rfloor}\).  For a
simplicial cone the bound is sharp.
\end{enumerate}
\end{theorem}
This theorem is quoted from Akian--Gaubert--Lemmens--Nussbaum and related work.  The survey gives
the concrete example on \(\R_+^3\)
\[
 f(x_1,x_2,x_3)=
 \begin{pmatrix}
 (3x_1\wedge x_2)\vee(3x_2\wedge x_3)\\
 (3x_1\wedge x_3)\vee(3x_3\wedge x_2)\\
 (3x_2\wedge x_1)\vee(3x_3\wedge x_1)
 \end{pmatrix},
\]
for which \((1,2,0)\) has period six.  Hence periods on \(\R_+^3\) include
\(1,2,3,4,6\) but not five.

\begin{theorem}[Theorem 4.5, quoted]
Let \((\Omega,\delta)\) be strictly convex and suppose no two-dimensional affine section
\(H\cap\Omega\) is an ellipsoid.  If a nonexpansive map \(f:\Omega\to\Omega\) has a fixed point,
then every orbit converges to a periodic orbit.  In fact there is a common \(q\ge1\) such that
\(f^{kq}(x)\) converges for every \(x\in\Omega\).
\end{theorem}

\subsection{Fixed-point-free maps and Denjoy--Wolff phenomena}
Calka's theorem, cited by the survey, says that if a Hilbert-nonexpansive map has no fixed point,
then every orbit's accumulation points lie on \(\partial\Omega\).  Define the attractor
\[
 A_f=\bigcup_{x\in\Omega}\omega(x;f)\subseteq\partial\Omega.
\]
The survey compares this problem with the classical disc theorem.

\begin{theorem}[Theorem 4.6, classical Denjoy--Wolff]
A fixed-point-free analytic self-map of the open unit disc has a unique boundary point
\(\eta\) such that \(f^k(x)\to\eta\) for every \(x\), uniformly on compact subsets.
\end{theorem}

\begin{theorem}[Theorem 4.7, Beardon, quoted]
If \((\Omega,\delta)\) is a strictly convex Hilbert geometry and
\(f:\Omega\to\Omega\) is fixed-point-free and nonexpansive, then there is a unique
\(\eta\in\partial\Omega\) with \(f^k(x)\to\eta\) for every \(x\), uniformly on compact subsets.
\end{theorem}
Strict convexity is essential: on the open two-simplex, Lins constructs fixed-point-free
nonexpansive maps whose total \(\omega\)-limit set is an arbitrary convex boundary subset.

The Karlsson--Nussbaum conjecture recorded by the survey asserts that for every fixed-point-free
Hilbert-nonexpansive map there should be a convex \(\Lambda\subseteq\partial\Omega\) containing
\(\omega(x;f)\) for every \(x\).  The paper lists the following partial results.

\begin{theorem}[Theorem 4.8, Lins, quoted]
The conjectured convex boundary set exists for every polytopal Hilbert geometry.
\end{theorem}
The survey explains the proof mechanism: Theorem 3.5 embeds the geometry in a finite-dimensional
normed space, allowing construction of a horofunction \(h\) with
\(h(f^k(x))\to-\infty\).  Such a horofunction need not exist for general Hilbert geometries.

\begin{theorem}[Theorem 4.9, Karlsson, quoted]
If some \(z\in\Omega\) has positive linear escape rate,
\[
 \lim_{k\to\infty}\frac{\delta(f^k(z),z)}k>0,
\]
then a convex \(\Lambda\subseteq\partial\Omega\) contains all \(\omega(x;f)\).
\end{theorem}

\begin{theorem}[Theorem 4.10, Karlsson--Noskov, quoted]
There is \(\eta\in\partial\Omega\) such that for every \(y\in A_f\), the Euclidean line segment
\([y,\eta]\) lies in \(\partial\Omega\).  Thus \(A_f\) is star-shaped in the stated boundary sense.
\end{theorem}

\begin{theorem}[Theorem 4.11, Nussbaum, quoted]
If some \(z\in\Omega\) satisfies
\[
 \liminf_{k\to\infty}\delta(f^{k+1}(z),f^k(z))=0,
\]
then a convex \(\Lambda\subseteq\partial\Omega\) contains all \(\omega(x;f)\).
\end{theorem}

The survey ends by emphasizing that these theorems do not settle the full Karlsson--Nussbaum
conjecture.  Their proofs belong to the cited papers; the survey's contribution here is the organized
connection between cone-induced Hilbert nonexpansiveness, finite-dimensional embeddings, and the
boundary-dynamics problem.

\section{Core theorem chain}
The core finite-dimensional statements are:
\begin{enumerate}[leftmargin=2em]
\item For positive vectors,
\[
  M(x/y)=\max_i x_i/y_i,
  \quad
  m(x/y)=\min_i x_i/y_i,
  \quad
  d(x,y)=\log(M/m).
\]
\item Order-preserving homogeneous degree-one maps are nonexpansive.
\item Strictly positive matrices have finite projective diameter via \eqref{eq:positive-matrix-diameter}.
\item Birkhoff's theorem turns finite diameter into a contraction factor \(\tanh(\Delta/4)<1\).
\item If a proof only needs asymptotic regularity or uniqueness of projective cluster points, it may be enough to combine finite diameter with monotone/nonexpansive arguments, postponing the exact contraction coefficient as a theorem seam.
\end{enumerate}

\begin{thebibliography}{99}
\bibitem{LemmensNussbaum2013} B. Lemmens and R. Nussbaum, \emph{Birkhoff's version of Hilbert's metric and its applications in analysis}, arXiv:1304.7921, 2013.
\bibitem{Birkhoff1957} G. Birkhoff, \emph{Extensions of Jentzsch's theorem}, Trans. Amer. Math. Soc. 85(1):219--227, 1957.
\bibitem{EvesonNussbaum1995} S. P. Eveson and R. D. Nussbaum, \emph{An elementary proof of the Birkhoff--Hopf theorem}, Math. Proc. Camb. Phil. Soc. 117:31--55, 1995.
\bibitem{FranklinLorenz1989} J. Franklin and J. Lorenz, \emph{On the scaling of multidimensional matrices}, Linear Algebra Appl. 114/115:717--735, 1989.
\bibitem{AkianEtAl2006} M. Akian, S. Gaubert, B. Lemmens, and R. D. Nussbaum, \emph{Iteration of order preserving subhomogeneous maps on a cone}, Math. Proc. Cambridge Philos. Soc. 140:157--176, 2006.
\bibitem{Beardon1997} A. F. Beardon, \emph{The dynamics of contractions}, Ergodic Theory Dynam. Systems 17:1257--1266, 1997.
\bibitem{Calka1984} A. Calka, \emph{On a condition under which isometries have bounded orbits}, Colloq. Math. 48:219--227, 1984.
\bibitem{Karlsson1999} A. Karlsson, \emph{Non-expanding maps and Busemann functions}, Ergodic Theory Dynam. Systems 21:1447--1457, 2001.
\bibitem{KarlssonNoskov2002} A. Karlsson and G. A. Noskov, \emph{The Hilbert metric and Gromov hyperbolicity}, Enseign. Math. 48:73--89, 2002.
\bibitem{Lins2008} B. Lins, results on nonexpansive maps and polytopal Hilbert geometries cited as references [42--43] in the survey.
\bibitem{LemmensScheutzow2005} B. Lemmens and M. Scheutzow, work on periods of nonexpansive maps cited as reference [38] in the survey.
\end{thebibliography}

\end{document}
